\documentclass[10pt,a4paper]{amsart}
\usepackage[margin=19mm]{geometry}
\usepackage{amsmath,amssymb,mathtools}
\usepackage[scr=rsfs,cal=euler]{mathalfa}
\usepackage{xfrac}
\usepackage[unicode,hidelinks,bookmarks=false]{hyperref}
\newcommand{\ie}{i.e.\ }
\newcommand{\cf}{cf.\ }
\newcommand{\resp}{respectively\ }
\newcommand{\Subsection}{\subsection}
\providecommand{\nobreakdash}{\nobreak\hbox{-}}
\setlength{\emergencystretch}{2em}
\multlinegap0pt
\usepackage{amssymb}
\usepackage{bm}
\newtheorem{proposition}{Proposition}
\title{Factorization and Atomic Decomposition in Hardy-Orlicz Spaces on the Upper Half-Plane with Applications to Hankel Operators}
\author{Jean-Marcel Tanoh Dje, Justin Feuto, Beno\^it F. Sehba}
\date{Source: arXiv:2608.19451v1 (CC BY 4.0)}
\begin{document}
\maketitle

\noindent\emph{Source and licence.} Factorization and Atomic Decomposition in Hardy-Orlicz Spaces on the Upper Half-Plane with Applications to Hankel Operators. Version arXiv:2608.19451v1 (CC BY 4.0), distributed by arXiv under the Creative Commons Attribution 4.0 International licence, http://creativecommons.org/licenses/by/4.0/, as recorded in the arXiv licence metadata for this exact version and shown on the abstract page https://arxiv.org/abs/2608.19451v1. Full licence notice: \texttt{licenses/arxiv-2608.19451-CC-BY-4.0.txt}.\par\smallskip

\noindent\textit{Editorial scope.} Counted unit: the article's proposition pro:mainqaq2aop (source0.tex lines 653--663). The article's complete original proof of it is delivered with it (source0.tex lines 665--858, one \texttt{proof} environment of 194 lines). No mathematics is written, repaired or re-derived: the statement and the proof are the article's own lines, in the article's own notation, and no retained line is substituted or re-flowed.  No further source-local context is needed: every cross-reference of every retained line resolves inside the two delivered spans.  Nothing was omitted from any retained span, so the package carries no omission record.  No retained line contains a citation, so the wrapper carries no bibliography and no bibliography entry.  No retained line contains a figure environment, an image-inclusion command or drawing code, so no image is delivered and the package declares no asset dependency.  Editorial formatting only, in addition: an AMS \texttt{amsart} preamble that restates the article's own declarations and macros used by the retained lines, namely the environments \texttt{proposition} on separate counters, together with the standard packages those lines need.  The retained environments print in this document as Proposition 1 in order of appearance, whereas the article prints the same items with the numbering of its own full document.  The complete native source of the article as distributed in the arXiv e-print for this exact version is bundled unchanged as \texttt{sources/208089/source0.tex}.
\begin{proposition}\label{pro:mainqaq2aop}
Let  $\Phi \in \mathscr{L}$ and $f \in \mathscr{S}'(\mathbb{R})$ be a bounded distribution. 
If $f \in H^{\Phi}(\mathbb{R})$, then for $t>0$, the function $P_{t}\ast f$ belongs to $H^{\Phi}(\mathbb{R})$ and 
\begin{equation}\label{eq:pmta2}
\|P_{t}\ast f\|_{H^{\Phi}(\mathbb{R})}^{lux} \leq C \|f\|_{H^{\Phi}(\mathbb{R})}^{lux},
\end{equation}
where $C$ is a constant independent of $f$ and $t$. Moreover,
\begin{equation}\label{eq:dqtaqa2}
\lim_{t \to 0}\|(P_{t}\ast f)-f\|_{H^{\Phi}(\mathbb{R})}^{lux}=0.
\end{equation}
\end{proposition}

\begin{proof}
Since $f$ is a bounded tempered distribution on $\mathbb{R}$, for every $s>0$ the convolution $P_s*f$ is well defined, bounded and continuous on $\mathbb{R}$. In particular, $(P_t)_{t>0}$ is an approximation of the identity. Hence, for all $x\in\mathbb{R}$ and all $s>0$,
\[
\lim_{t\to 0} P_t*(P_s*f)(x)=P_s*f(x).
\]
Using the semigroup property $P_t*(P_s*f)=P_{s+t}*f$, we obtain
\begin{equation}\label{eq:pointwise_limit}
\lim_{t\to 0} P_{s+t}*f(x)=P_s*f(x),
\qquad \forall s>0,\ \forall x\in\mathbb{R}.
\end{equation}
For $t>0$ and $s>0$,
\[
|P_s*(P_t*f)(x)|
=
|P_{s+t}*f(x)|
\le
\sup_{u>0}|P_u*f(x)|
=
f_P^*(x).
\]
Thus
\[
(P_t*f)_P^*(x)\le f_P^*(x).
\]
By the definition of the Hardy–Orlicz norm,
\[
\|P_t*f\|_{H^\Phi(\mathbb{R})}^{lux}
\approx
\|(P_t*f)_P^*\|_{L^\Phi}^{lux}
\le
\|f_P^*\|_{L^\Phi}^{lux}
\approx
\|f\|_{H^\Phi(\mathbb{R})}^{lux}.
\]
Let
\[
u(s,x):=P_s*f(x), \qquad \forall~s>0,\ \forall~ x\in\mathbb{R}.
\]
It is well known that $u(s,\cdot)\longrightarrow f$
in  $\mathscr S'(\mathbb{R})$  as  $s\to 0$.
Hence, for every $\varepsilon>0$ and every $\varphi\in\mathscr S(\mathbb{R})$, there exists
$k_{\varepsilon}(\varphi)>0$ such that
\[
|\langle u(s,\cdot)-f,\varphi\rangle|
<
\varepsilon,
\qquad \forall~ 0<s<k_{\varepsilon}(\varphi).
\]
Let $x\in\mathbb{R}$.  
Choose $\psi\in\mathscr S(\mathbb{R})$ such that
$\int_{\mathbb R}\psi(y)\,dy=1$,
and define
\[
\psi_\eta(y)=\frac1\eta\,\psi\!\left(\frac{x-y}{\eta}\right),
\qquad \forall~ \eta>0 .
\]
Then $(\psi_\eta)_{\eta>0}$ is an approximation of the Dirac mass at $x$, that is $\psi_\eta\longrightarrow \delta_x$ in  $\mathscr S'(\mathbb{R})$  as  $\eta\to 0$.
Applying the previous estimate with $\varphi=\psi_\eta$, we obtain that for every
$\varepsilon>0$ there exists $k_{\varepsilon}(\psi_\eta)>0$ such that
\[
|\langle u(s,\cdot)-f,\psi_\eta\rangle|
<
\varepsilon,
\qquad \forall~ 0<s<k_{\varepsilon}(\psi_\eta).
\]
In particular, for all $0<s,t<k_{\varepsilon}(\psi_\eta)/2$,
\[
|\langle u(s+t,\cdot)-u(s,\cdot),\psi_\eta\rangle|
\le
|\langle u(s+t,\cdot)-f,\psi_\eta\rangle|
+
|\langle u(s,\cdot)-f,\psi_\eta\rangle|
<
2\varepsilon .
\]
Since
\[
\langle u(s+t,\cdot)-u(s,\cdot),\psi_\eta\rangle
=
\int_{\mathbb R}
\bigl(u(s+t,y)-u(s,y)\bigr)\psi_\eta(y)\,dy,
\]
and since $u(s,\cdot)$ is continuous, the family $(\psi_\eta)$ converges to $\delta_x$,
so that
\[
\lim_{\eta\to 0}
\langle u(s+t,\cdot)-u(s,\cdot),\psi_\eta\rangle
=
u(s+t,x)-u(s,x).
\]
Consequently, letting $\eta\to 0$ in the previous inequality yields
\[
|u(s+t,x)-u(s,x)|
\le
2\varepsilon,
\qquad \forall~ 0<s,t<k_{\varepsilon}(x)/2,
\]
for almost every $x\in\mathbb R$. Therefore
\begin{equation}\label{eq:small_s}
\sup_{0<s<k_\varepsilon(x)/2}
|u(s+t,x)-u(s,x)|
\le
2\varepsilon,
\qquad \forall~ 0<t<k_\varepsilon(x)/2 .
\end{equation}
On the other hand, since the Poisson kernel is smooth for $s>0$, the mapping
$s\mapsto u(s,x)$ is $\mathscr C^1$ and
\[
\partial_s u(s,x)
=
(\partial_s P_s)*f(x).
\]
It is well known that
\[
|\partial_s P_s(y)|
\lesssim
\frac1s P_s(y).
\]
Hence
\[
|\partial_s u(s,x)|
\lesssim
\frac1s f_P^*(x).
\]
Using the fundamental theorem of calculus,
\[
u(s+t,x)-u(s,x)
=
\int_s^{s+t}\partial_r u(r,x)\,dr.
\]
Therefore
\[
|u(s+t,x)-u(s,x)|
\lesssim
f_P^*(x)\int_s^{s+t}\frac{dr}{r}
\lesssim
\frac{t}{s}f_P^*(x).
\]
Thus
\begin{equation}\label{eq:large_s}
\sup_{s\ge k_\varepsilon(x)/2}
|u(s+t,x)-u(s,x)|
\lesssim
\frac{t}{k_\varepsilon(x)}f_P^*(x).
\end{equation}
Combining \eqref{eq:small_s} and \eqref{eq:large_s}, we obtain
\[
(P_t*f-f)_P^*(x)
=
\sup_{s>0}|u(s+t,x)-u(s,x)|
\lesssim
2\varepsilon
+
\frac{t}{k_\varepsilon(x)}f_P^*(x), 
\qquad \forall~ 0<s,t<k_{\varepsilon}(x)/2,
\]
for almost every $x\in\mathbb R$. Letting $t\to 0$ gives
\[
(P_t*f-f)_P^*(x)\longrightarrow0
\quad \text{for almost every } x\in\mathbb R .
\]
Moreover,
\[
(P_t*f-f)_P^*(x)
\le
2f_P^*(x).
\]
Since $\Phi\in\mathscr L$, there exists $K>0$ such that
\[
\Phi(2u)\le K\Phi(u).
\]
Thus
\[
\Phi((P_t*f-f)_P^*(x))
\le
K\Phi(f_P^*(x)).
\]
Because $f\in H^\Phi(\mathbb R)$, we have
$\Phi(f_P^*)\in L^1(\mathbb R)$.
Hence, by the dominated convergence theorem,
\[
\int_{\mathbb R}
\Phi\!\left(\varepsilon (P_t*f-f)_P^*\right)dx
\longrightarrow0
\quad (t\to0).
\]
By the definition of the Luxemburg norm,
\[
\|P_t*f-f\|_{H^\Phi(\mathbb R)}^{lux}
\longrightarrow0 .
\]

This completes the proof.
\end{proof}

\end{document}
